% 7.1：沿用问题二的全过程模型，新增全域达标判据与未知终止时间。
\subsection{模型建立}
\label{sec:p3-model}

问题三在问题二模型上增加全域含水率约束，
要求各处均低于$0.15\cunit$，由此确定未知的临界烘干时间。

\subsubsection{全过程模型的沿用与延续}

沿用第\ref{sec:p2-model}节的全过程模型，以及式\eqref{eq:p1-environment}的插值和前文规定的恒定环境延拓，
将问题二在$3\ \mathrm{h}$时的两场状态继续向后推进，
总时长仍从原始初态$t=0$起算。

\subsubsection{全域达标约束与临界时间}

令$C_*=0.15\cunit$。全域达标要求每个径向位置均满足$C(r,t)<C_*$，
等价于$C_{\max}(t)<C_*$，其中全域最大含水率及临界时间定义为
\begin{align}
  C_{\max}(t)&=\max_{0\le r\le R_0}C(r,t),
  \label{eq:p3-maximum}\\
  t_*&=\inf\left\{t\ge0:\ C_{\max}(t)<C_*\right\}.
  \label{eq:p3-endpoint}
\end{align}
当$C_{\max}(t)$连续向下穿越阈值时，临界处$C_{\max}(t_*)=C_*$，
其后才严格达标；判断须基于未舍入的场量。
烘干结束由含水率最高的位置控制，表面或截面平均含水率不能替代全域判据。
中心能否代表全域最大值，须在求解中逐步核验。
